\subsection{Theoretical Background: The Harmonic Oscillator \quad / \quad \textthai{พื้นฐานทางทฤษฎี: ตัวสั่นพ้องแบบฮาร์มอนิก}}
\begin{paracol}{2}
At the core of many physical phenomena—from the vibration of atoms in a crystal to the oscillations of a guitar string—lies the simple harmonic oscillator\index{Simple Harmonic Oscillator@ตัวสั่นพ้องแบบฮาร์มอนิกอย่างง่าย (Simple Harmonic Oscillator)} (SHO). It is the most fundamental model for any system close to a stable equilibrium point.

Mathematically, a simple harmonic oscillator is described by Hooke's Law\index{Hooke's Law@กฎของฮุค (Hooke's Law)}, where the restoring force $F$ is proportional to the displacement $x$:
\begin{equation}
    F = -kx
\end{equation}
Applying \textbf{Newton's Second Law}\index{Newton's Second Law@กฎข้อที่สองของนิวตัน (Newton's Second Law)} ($F = ma = m \frac{d^2x}{dt^2}$), we arrive at the second-order linear \textbf{ordinary differential equation (ODE)}\index{ODE@สมการเชิงอนุพันธ์สามัญ (ODE)}:
\begin{equation}
    m \frac{d^2x}{dt^2} + kx = 0
\end{equation}
Dividing by $m$ and defining the natural frequency $\omega_0^2 = k/m$, we get the standard form:
\begin{equation}
    \frac{d^2x}{dt^2} + \omega_0^2 x = 0
\end{equation}

The general solution to this equation is a combination of sines and cosines:
\begin{equation}
    x(t) = A \cos(\omega_0 t + \phi)
\end{equation}

\switchcolumn

\begin{thai}
ที่แกนกลางของปรากฏการณ์ทางกายภาพหลายอย่าง ตั้งแต่การสั่นสะเวือนของอะตอมในผลึกไปจนถึงการแกว่งของสายกีตาร์ คือตัวสั่นพ้องแบบฮาร์มอนิกอย่างง่าย (SHO) ซึ่งเป็นแบบจำลองพื้นฐานที่สุดสำหรับระบบใดก็ตามที่อยู่ใกล้จุดสมดุลที่เสถียร

ในทางคณิตศาสตร์ ตัวสั่นพ้องแบบฮาร์มอนิกอย่างง่ายอธิบายได้ด้วยกฎของฮุค (Hooke's Law) ซึ่งแรงดึงกลับ $F$ จะแปรผันตรงกับระยะขจัด $x$:
\begin{equation}
    F = -kx
\end{equation}
เมื่อใช้กฎข้อที่สองของนิวตัน ($F = ma = m \frac{d^2x}{dt^2}$) เราจะได้\textbf{สมการเชิงอนุพันธ์สามัญ (ODE)}\index{ODE@สมการเชิงอนุพันธ์สามัญ (ODE)}เชิงเส้นอันดับสอง:
\begin{equation}
    m \frac{d^2x}{dt^2} + kx = 0
\end{equation}
เมื่อหารด้วย $m$ และกำหนดความถี่ธรรมชาติ $\omega_0^2 = k/m$ เราจะได้รูปแบบมาตรฐาน:
\begin{equation}
    \frac{d^2x}{dt^2} + \omega_0^2 x = 0
\end{equation}

คำตอบทั่วไปของสมการนี้คือการรวมกันของฟังก์ชันไซน์และคอสไซน์:
\begin{equation}
    x(t) = A \cos(\omega_0 t + \phi)
\end{equation}
\end{thai}
\end{paracol}
